\documentclass[8pt, landscape, a4paper]{scrartcl}
\usepackage[utf8]{inputenc}
\usepackage[margin=1cm]{geometry}
\usepackage{amsmath, amssymb}
\usepackage{multicol}
\usepackage{enumitem}
\usepackage{titlesec}
\usepackage{xcolor}
\usepackage{array}

% --- Style Definitions ---
\setlength{\parindent}{0pt}
\setlength{\parskip}{0pt}
\setlist{nosep} % Compact lists

% Compact section spacing
\titlespacing*{\section}{0pt}{4pt}{2pt}
\titlespacing*{\subsection}{0pt}{3pt}{1pt}
\titlespacing*{\subsubsection}{4pt}{0pt}{0pt}

% Section styling
\titleformat{\section}
  {\normalfont\large\bfseries\color{blue}}
  {\thesection}{1em}{}
\titleformat{\subsection}
  {\normalfont\normalsize\bfseries}
  {\thesubsection}{1em}{}
\titleformat{\subsubsection}
    {\normalfont\small\itshape}
    {\thesubsubsection}{1em}{}
% -------------------------

\begin{document}

\begin{center}
    {\Huge \textbf{MA1521 Cheatsheet}}
\end{center}

\begin{multicols*}{4}

\section*{Ch 0: Real Numbers \& Functions}
\subsection*{Properties}
\begin{itemize}
    \item \textbf{Triangle Inequality:} $|x + y| \le |x| + |y|$
    \item \textbf{Absolute Value:}
    \begin{itemize}
        \item $|-x| = |x|$
        \item $|xy| = |x||y|$
        \item $\sqrt{x^2} = |x|$
        \item $|x| < r \iff -r < x < r$ (for $r > 0$)
    \end{itemize}
\end{itemize}

\section*{Ch 1: Limits and Continuity}
\subsection*{Theorems \& Laws}
\begin{itemize}
    \item \textbf{Limit Laws:} For sums, differences, products, and quotients (denom $\neq 0$), the limit of the operation is the operation of the limits.
    \item \textbf{Squeeze (Sandwich) Theorem:} If $g(x) \le f(x) \le h(x)$ near $c$ and $\lim_{x \to c} g(x) = \lim_{x \to c} h(x) = L$, then $\lim_{x \to c} f(x) = L$.
    \item \textbf{Intermediate Value Theorem (IVT):} If $f$ is continuous on $[a, b]$ and $k$ is between $f(a)$ and $f(b)$, there exists $c \in [a, b]$ such that $f(c) = k$.
\end{itemize}

\subsection*{Special Trig Limits}
\begin{itemize}
    \item \textbf{Small Angle Approximation:}
    \[ \lim_{x \to 0} \frac{\sin x}{x} = 1 \quad \text{and} \quad \lim_{x \to 0} \frac{1 - \cos x}{x} = 0 \]
    \item \textbf{Note:} $\lim_{n \to \infty} \frac{\sin n}{n} = 0$ (by Squeeze Thm).
    \item \textit{Application:} $\lim_{n \to \infty} n \sin\left(\frac{k}{n}\right) = k$.
\end{itemize}

\subsection*{Continuity}
$f$ is continuous at $c$ if:
\[ \lim_{x \to c} f(x) = f(c) \]

\section*{Ch 2: Derivatives}
\subsection*{Definition}
\[ f'(x) = \lim_{h \to 0} \frac{f(x+h) - f(x)}{h} \]

\subsection*{Table of Derivatives}
\noindent
\begin{tabular}{@{}ll@{}}
    $(x^n)' = nx^{n-1}$   & $(\cot x)' = -\csc^2 x$ \\
    $(e^x)' = e^x$        & $(\sec x)' = \sec x \tan x$ \\
    $(\ln x)' = 1/x$      & $(\csc x)' = -\csc x \cot x$ \\
    $(\sin x)' = \cos x$  & $(\sin^{-1} x)' = \frac{1}{\sqrt{1-x^2}}$ \\
    $(\cos x)' = -\sin x$ & $(\tan^{-1} x)' = \frac{1}{1+x^2}$ \\
    $(\tan x)' = \sec^2 x$& $(\sec^{-1} x)' = \frac{1}{|x|\sqrt{x^2-1}}$ \\
\end{tabular}

\subsection*{Differentiation Rules}
\begin{itemize}
    \item \textbf{Product Rule:}
    \[ \frac{d}{dx}(uv) = u\frac{dv}{dx} + v\frac{du}{dx} \]
    \item \textbf{Quotient Rule:}
    \[ \frac{d}{dx}\left(\frac{u}{v}\right) = \frac{v\frac{du}{dx} - u\frac{dv}{dx}}{v^2} \]
    \item \textbf{Chain Rule:} If $y = f(u)$ and $u = g(x)$:
    \[ \frac{dy}{dx} = \frac{dy}{du} \cdot \frac{du}{dx} \]
\end{itemize}

\subsection*{Parametric Differentiation}
For $x=f(t), y=g(t)$:
\[ \frac{dy}{dx} = \frac{g'(t)}{f'(t)} \quad \text{and} \quad \frac{d^2y}{dx^2} = \frac{\frac{d}{dt}(y')}{\frac{dx}{dt}} \]

\section*{Ch 3: Applications of Differentiation}
\subsection*{Lines \& Approximations}
\begin{itemize}
    \item \textbf{Tangent Line:} $y - y_0 = m(x - x_0)$ where $m = f'(x_0)$.
    \item \textbf{Normal Line:} $y - y_0 = -\frac{1}{m}(x - x_0)$.
    \item \textbf{Linear Approximation (Linearization):}
    \[ f(x) \approx L(x) = f(a) + f'(a)(x-a) \]
    \item \textbf{Newton's Method (Root Finding):}
    \[ x_{n+1} = x_n - \frac{f(x_n)}{f'(x_n)} \]
\end{itemize}

\subsection*{Theorems}
\begin{itemize}
    \item \textbf{Rolle's Theorem:} If $f$ is continuous on $[a, b]$, differentiable on $(a, b)$, and $f(a) = f(b)$, there is a $c \in (a, b)$ where $f'(c) = 0$.
    \item \textbf{Mean Value Theorem (MVT):} Under the same conditions, there is a $c \in (a, b)$ where:
    \[ f'(c) = \frac{f(b) - f(a)}{b - a} \]
    \item \textbf{L'Hôpital's Rule:} For forms $\frac{0}{0}$ or $\frac{\infty}{\infty}$:
    \[ \lim_{x \to c} \frac{f(x)}{g(x)} = \lim_{x \to c} \frac{f'(x)}{g'(x)} \]
\end{itemize}

\subsection*{Extrema Tests}
\begin{itemize}
    \item \textbf{First Derivative Test:}
    \begin{itemize}
        \item Change $+$ to $-$ $\implies$ Local Max
        \item Change $-$ to $+$ $\implies$ Local Min
    \end{itemize}
    \item \textbf{Second Derivative Test:}
    \begin{itemize}
        \item $f'(c)=0, f''(c) < 0 \implies$ Local Max
        \item $f'(c)=0, f''(c) > 0 \implies$ Local Min
    \end{itemize}
\end{itemize}

\section*{Ch 4: Integrals}
\subsection*{Riemann Sum Definition}
The definite integral $\int_{a}^{b} f(x) \, dx$ is defined as:
\[ \lim_{n\to\infty} \left\{ \sum_{k=1}^{n} \left(\frac{b-a}{n}\right) f\left(a + k\left(\frac{b-a}{n}\right)\right) \right\} \]

\subsection*{Fundamental Theorem of Calculus}
\begin{itemize}
    \item \textbf{Part 1:} $\int_{a}^{b} f(x) \, dx = F(b) - F(a)$ ($F$ is antiderivative of $f$).
    \item \textbf{Part 2:} $\frac{d}{dx} \int_{a}^{x} f(t) \, dt = f(x)$.
\end{itemize}

\subsection*{Common \& Trig Integrals}
\begin{itemize}
    \item $\int \ln x \, dx = x \ln x - x + C$
    \item $\int \tan x \, dx = \ln |\sec x| + C$
    \item $\int \cot x \, dx = \ln |\sin x| + C$
    \item $\int \sec x \, dx = \ln |\sec x + \tan x| + C$
    \item $\int \csc x \, dx = \ln |\csc x - \cot x| + C$
\end{itemize}

\subsection*{Techniques}
\begin{itemize}
    \item \textbf{Integration by Parts:}
    \[ \int u \, dv = uv - \int v \, du \]
    \item \textbf{Substitution:}
    \[ \int f(g(x))g'(x) \, dx = \int f(u) \, du \]
\end{itemize}

\subsection*{Tangent Half-Angle ($z=\tan\left(\frac{\theta}{2}\right)$)}
Used for rational functions of $\sin\theta, \cos\theta$:
\begin{itemize}
    \item $\sin\theta = \frac{2z}{1+z^2}, \quad \cos\theta = \frac{1-z^2}{1+z^2}$
    \item $d\theta = \frac{2}{1+z^2} \,dz$
    \item \textbf{Limits:} $0 \to 0, \quad \frac{\pi}{2} \to 1, \quad \pi \to \infty$
\end{itemize}

\section*{Ch 5: Applications of Integration}
\subsection*{Area Between Curves}
\begin{itemize}
    \item \textbf{Vertical Strips:} $A = \int_{a}^{b} |f(x) - g(x)| \, dx$
    \item \textbf{Horizontal Strips:} $A = \int_{c}^{d} |f(y) - g(y)| \, dy$
\end{itemize}

\subsection*{Volume of Revolution}
\begin{itemize}
    \item \textbf{Disk Method (about x-axis):}
    \[ V = \pi \int_{a}^{b} {[f(x)]}^2 \, dx \]
    \item \textbf{Washer Method (about x-axis):}
    \[ V = \pi \int_{a}^{b} \left({[f(x)]}^2 - {[g(x)]}^2\right) \, dx \]
    \item \textbf{Shell Method (about y-axis):}
    \[ V = 2\pi \int_{a}^{b} x|f(x)| \, dx \]
\end{itemize}

\subsection*{Surface Area of Revolution}
\begin{itemize}
    \item \textbf{About the x-axis:}
    For $y = f(x)$ from $a$ to $b$:
    \[ S = 2\pi \int_{a}^{b} f(x) \sqrt{1 + {[f'(x)]}^2} \, dx \]
    \item \textbf{About the y-axis:}
    For $x = g(y)$ from $c$ to $d$:
    \[ S = 2\pi \int_{c}^{d} g(y) \sqrt{1 + {[g'(y)]}^2} \, dy \]
    \item \textbf{For parametric equations (about x-axis):}
    \[ S = 2\pi \int_c^d y(t) \sqrt{{\left(\frac{dx}{dt}\right)}^2 + {\left(\frac{dy}{dt}\right)}^2} \, dt \]
\end{itemize}

\subsection*{Arc Length}
\begin{itemize}
    \item For $y = f(x)$ from $a$ to $b$:
    \[ L = \int_{a}^{b} \sqrt{1 + {[f'(x)]}^2} \, dx \]
    \item For $x = g(y)$ from $c$ to $d$:
    \[ L = \int_{c}^{d} \sqrt{1 + {[g'(y)]}^2} \, dy \]
    \item For parametric equations:
    \[L = \int_c^d \sqrt{{\left(\frac{dx}{dt}\right)}^2 + {\left(\frac{dy}{dt}\right)}^2} \,dt\]
\end{itemize}

\section*{Ch 6: Infinite/Convergent Series}

\subsection*{Absolute vs Conditional Convergence}
\begin{itemize}
    \item \textbf{Absolute Convergence:} $\sum|a_n|$ converges
    \item \textbf{Conditional Convergence:} $\sum a_n$ converges, but $\sum |a_n|$ diverges
\end{itemize}

\subsection*{Sum of Geometric Series}
\[
    S_n = \sum_{i=1}^{n} a_i \cdot r^{i-1} = a_1\left(\frac{1-r^n}{1-r}\right)
\]
\begin{itemize}
    \item Diverges if $|r| \geq 1$
    \item Converges to $\frac{a}{1-r}$ if $|r| < 1$
\end{itemize}

\subsection*{Series}
    \subsubsection*{p-Series}
    \[
        \sum_{n=1}^{\infty} \frac{1}{n^p}
    \]
    
    \begin{itemize}
        \item Converges if $p > 1$
        \item Diverges if $p \leq 1$
    \end{itemize}
    \subsubsection*{Harmonic Series}
    \[
        \sum_{n=1}^{\infty} \frac{1}{n}
    \]
    \begin{itemize}
        \item \textbf{Never Converges}
    \end{itemize}
\begin{itemize}
    \item Specific case of \textbf{p-series} when \textit{p = 1}
\end{itemize}

\subsection*{Test for Convergence}
    \subsubsection*{Integral Test}
    \begin{itemize}
        \item For a series $\sum a_n$ where $a_n = f(n)$ and $f(x)$ are continuous, positive, and decreasing for $x \geq 1$:
        \item Converges \textbf{if and only if:}
    \end{itemize}
    \[
        \int_{1}^{\infty}f(x)\,dx \text{ \textbf{converges}}
    \]

    \subsubsection*{Comparison Test}
    \begin{itemize}
        \item {Suppose $0 \leq a_n \leq b_n$}
        \begin{itemize}
            \item {If $\sum b_n$ converges, then $\sum a_n$ converges}
            \item {If $\sum a_n$ diverges, then $\sum b_n$ diverges}
        \end{itemize}
    \end{itemize}

    \subsection*{Ratio Test}
    \[
        L = \lim_{n\to\infty} \left|\frac{a_{n+1}}{a_n}\right|
    \]
    \begin{itemize}
        \item If $L < 1$: \textbf{Absolutely Convergent}
        \item If $L > 1$: \textbf{Divergent}
        \item If $L = 1$: \textbf{Inconclusive}
    \end{itemize}

    \subsection*{Root Test}
    \[
        L = \lim_{n \to \infty} \sqrt[n]{|a_n|}
    \]
    \begin{itemize}
        \item Conclusions are identical to Ratio Test
    \end{itemize}

    \subsection*{Alternating Series Test}
    \[
        \sum {(-1)}^{n-1}b_n \text{ with } b_n > 0
    \]
    \begin{itemize}
        \item If $b_n$ is decreasing $(b_{n+1} \leq b_n)$ and $\lim_{n\to\infty}b_n = 0$, then the series converges.
    \end{itemize}

    \subsection*{Power Series}
    \[
        \sum_{n = 0}^{\infty} c_n \cdot {(x-a)}^n \text{, centered at } a
    \]
    \begin{itemize}
        \item Radius of Convergence
        \begin{itemize} 
            \item The series converges abs for $|x-a| < R$ and diverges for $|x-a| > R$
            \item $R = \frac{1}{L} \implies L$ is the limit from Ratio or Root test applied to coeffs
        \end{itemize}
        \item \textbf{Take note:} Within radius of convergence, can be differentiated and integrated term-by-term
    \end{itemize}

\subsection*{Taylor and Maclaurin Series}
\subsubsection*{Taylor Series}
\[
    f(x)=\sum_{n=0}^\infty \frac{f^{(n)}(a)}{n!}{(x-a)}^n
\]
\begin{itemize}
    \item Representation of a function $f(x)$ as a power series centred at $x=a$
\end{itemize}

\subsubsection*{Maclaurin Series}
\begin{itemize}
    \item Taylor Series centred at $a=0$
\end{itemize}

\subsubsection*{Important Maclaurin Series}
\begin{itemize}
    \item $e^x = \sum_{n=0}^{\infty} \frac{x^n}{n!} = 1 + x + \frac{x^2}{2!} + \frac{x^3}{3!} + \dots$
    \item $(1+x)^k = 1 + kx + \frac{k(k-1)}{2!}x^2 + \frac{k(k-1)(k-2)}{3!}x^3 + \dots$
    \item $ln(1+x) = x - \frac{x^2}{2} + \frac{x^3}{3} - \frac{x^4}{4} + \dots$ (converges for $-1 < x \leq 1$)
    \item $\sin x = \sum_{n=0}^{\infty} \frac{(-1)^n x^{2n+1}}{(2n+1)!} = x - \frac{x^3}{3!} + \frac{x^5}{5!} - \dots$
    \item $\cos x = \sum_{n=0}^{\infty} \frac{(-1)^n x^{2n}}{(2n)!} = 1 - \frac{x^2}{2!} + \frac{x^4}{4!} - \dots$
    \item $\tan^{-1} x = \sum_{n=0}^{\infty} \frac{(-1)^n x^{2n+1}}{2n+1} = x - \frac{x^3}{3} + \frac{x^5}{5} - \dots$ (converges for $|x| \leq 1$)
\end{itemize}

\section*{Ch 7: Vectors and Geometry of Space}
\subsection*{3D Coordinates \& Vectors}
\begin{itemize}
    \item \textbf{Distance:} \[|P_1 P_2| = \sqrt{{(x_2-x_1)}^2 + {(y_2-y_1)}^2 + {(z_2-z_1)}^2}\]
    \item \textbf{Equation of Sphere:} \[{(x-h)}^2 + {(y-k)}^2 + {(z-l)}^2 = r^2\]
    \item \textbf{Length:} $\lVert\mathbf{a}\rVert = \sqrt{{a_1}^2 + {a_2}^2 + {a_3}^2}$
    \item \textbf{Unit Vector:} $\mathbf{u} = \frac{\mathbf{a}}{\lVert\mathbf{a}\rVert}$
\end{itemize}

\subsection*{The Dot Product}
\begin{itemize}
    \item $\mathbf{a} \cdot \mathbf{b} = a_1 b_1 + a_2 b_2 + a_3 b_3 = \lVert\mathbf{a}\rVert \lVert\mathbf{b}\rVert \cos \theta$
    \item $\mathbf{a} \perp \mathbf{b} \iff \mathbf{a} \cdot \mathbf{b} = 0$
\end{itemize}

\subsection*{Projections}
\begin{itemize}
    \item \textbf{Scalar Projection (comp):} \[\text{comp}_\mathbf{a} \mathbf{b} = \frac{\mathbf{a} \cdot \mathbf{b}}{\lVert\mathbf{a}\rVert}\]
    \item \textbf{Vector Projection (proj):} \[\text{proj}_\mathbf{a} \mathbf{b} = \left(\frac{\mathbf{a} \cdot \mathbf{b}}{\lVert\mathbf{a}\rVert^2}\right) \mathbf{a}\]
    \item \textbf{Point to Plane Distance:} \[D = \frac{|ax_0+by_0+cz_0-d|}{\sqrt{a^2+b^2+c^2}}\]
\end{itemize}

\subsection*{The Cross Product}
\begin{itemize}
    \item \textbf{Definition:} 
            \[
            \mathbf{a} \times \mathbf{b} = 
            \begin{vmatrix} 
            \mathbf{i} & \mathbf{j} & \mathbf{k} \\ 
            a_1 & a_2 & a_3 \\ 
            b_1 & b_2 & b_3 
            \end{vmatrix}
            \]
    \item \textbf{Magnitude:} $\lVert\mathbf{a} \times \mathbf{b}\rVert = \lVert\mathbf{a}\rVert \lVert\mathbf{b}\rVert \sin \theta$
    \item $\mathbf{a} \times \mathbf{b}$ is orthogonal to both $\mathbf{a}$ and $\mathbf{b}$.
    \item \textbf{Point Q to Line PR Dist:} $d = \frac{\lVert\vec{PQ} \times \vec{PR}\rVert}{\lVert\vec{PR}\rVert}$
\end{itemize}

\subsection*{Lines \& Planes}
\begin{itemize}
    \item \textbf{Line (Vector):} $\mathbf{r} = \mathbf{r}_0 + t\mathbf{v}$
    \item \textbf{Line (Parametric):} $x = x_0+at$, $y = y_0+bt$, $z = z_0+ct$
    \item \textbf{Plane (Vector):} $\mathbf{n} \cdot (\mathbf{r} - \mathbf{r}_0) = 0$
    \item \textbf{Plane (Scalar):} $A(x-x_0) + B(y-y_0) + C(z-z_0) = 0$
\end{itemize}

\section*{Ch 8: Functions of Several Variables}
    \subsection*{Arc Length Formula}
    \begin{itemize}
        \item Let C be: $r(t) = \langle f(t), g(t), h(t) \rangle$, $a\leq t \leq b$
        \item $f'$, $g'$ and $h'$ are continuous
    \end{itemize}
    \[s = \int_a^b \sqrt{{f'(t)}^2 + {g'(t)}^2 + {h'(t)}^2} \,dt\]

    \subsection*{Partial Derivatives}
    \[f_x(x,y) = \lim_{h \to 0} \frac{f(x+ h, y) - f(x, y)}{h}\]
    \[f_y(x,y) = \lim_{h \to 0} \frac{f(x, y +h) - f(x,y)}{h}\]

    \subsection*{Clairaut's Theorem}
    If f is defined on a disk D that contains (a,b), and $f_{xy}$ and $f_{yx}$ are both continuous on D, then
    \[f_{xy}(a,b) = f_{yx}(a,b)\]

    \subsection*{Tangent Plane \& Linear Approx}
    \begin{itemize}
        \item \textbf{Tangent Plane:} $z = f(a,b) + f_x (a,b)(x-a) + f_y(a,b)(y - b)$
        \item \textbf{Linear Approximation:} $f(x,y) \approx f(a,b) + f_x(a,b)dx + f_y(a,b)dy$
    \end{itemize}
    
    \subsection*{Second Partials Test}
    Let $D = f_{xx}(a,b)f_{yy}(a,b) - [f_{xy}(a,b)]^2$:
    \begin{itemize}
        \item $D > 0, f_{xx}(a,b) > 0 \implies$ Local Minimum
        \item $D > 0, f_{xx}(a,b) < 0 \implies$ Local Maximum
        \item $D < 0 \implies$ Saddle Point
        \item $D = 0 \implies$ Test is Inconclusive
    \end{itemize}

    \subsection*{Chain Rule}
    \begin{itemize}
        \item Case 1: $z = f(x,y)$, where $x = g(t)$ and $y = h(t)$ (all are differentiable by t)
        \[\frac{dz}{dt} = \frac{\partial f}{\partial x} \frac{\partial x}{\partial t} + \frac{\partial f}{\partial y}\frac{\partial y}{\partial t}\]
        \item Case 2: $z = f(x,y)$, where $x = g(s,t)$ and $y = h(s, t)$ (all are differentiable by s and t)
        \[\frac{\partial z}{\partial s} = \frac{\partial f}{\partial x} \frac{\partial x}{\partial s} + \frac{\partial f}{\partial y}\frac{\partial y}{\partial s}\]
        \[\frac{\partial z}{\partial t} = \frac{\partial f}{\partial x} \frac{\partial x}{\partial t} + \frac{\partial f}{\partial y}\frac{\partial y}{\partial t}\]
        \item General Case, where u is a differentiable function of n variables $x_1, \dots, x_n$ and each $x_j$ is a differentiable function of m variables $t_1, \dots, t_m$. Then u is a function of $t_1, \dots, t_m$:
        \[\frac{\partial u}{\partial t_i} = \frac{\partial u}{\partial x_1} \frac{\partial x_1}{\partial t_i} + \frac{\partial u}{\partial x_2}\frac{\partial x_2}{\partial t_i}+\cdots+\frac{\partial u}{\partial x_n} \frac{\partial x_n}{\partial t_i}\]

        \textbf{for each $i = 1, \dots, m$}
    \end{itemize}
\section*{Trigonometric Identities}
\subsection*{Pythagorean}
\begin{itemize}
    \item $\sin^2 x + \cos^2 x = 1$
    \item $\sec^2 x - 1 = \tan^2 x$
\end{itemize}

\subsection*{Double Angle}
\begin{itemize}
    \item $\sin 2x = 2 \sin x \cos x$
    \item $\cos 2x = 2 \cos^2 x - 1 = 1 - 2 \sin^2 x$
\end{itemize}

\subsection*{Power-Reduction (Half-Angle)}
\[
    \cos^2 x = \frac{1 + \cos 2x}{2}
\]
\[
    \sin^2 x = \frac{1 - \cos 2x}{2}
\]

\subsection*{Product-to-Sum}
\[
    \sin A \cos B = \frac{1}{2}\left[\sin(A+B) + \sin(A-B)\right]
\]
\[
    \cos A \cos B = \frac{1}{2}\left[\cos(A+B) + \cos(A-B)\right]
\]

\subsection*{Special Inverse \& Hyperbolic Integrals}
\begin{itemize}
    \item $\int \frac{1}{\sqrt{a^2-x^2}} \, dx = \sin^{-1}\left(\frac{x}{a}\right) + C$
    \item $\int \frac{1}{a^2+x^2} \, dx = \frac{1}{a}\tan^{-1}\left(\frac{x}{a}\right) + C$
    \item $\int \frac{1}{x\sqrt{x^2-a^2}} \, dx = \frac{1}{a}\sec^{-1}\left|\frac{x}{a}\right| + C$
    \item $\int \frac{1}{\sqrt{x^2+a^2}} \, dx = \ln|x + \sqrt{x^2+a^2}| + C$
    \item $\int \frac{1}{\sqrt{x^2-a^2}} \, dx = \ln|x + \sqrt{x^2-a^2}| + C$
\end{itemize}

\subsection*{Trigonometric Substitution}
Used for integrals containing radicals of quadratic forms:
\begin{itemize}
    \item \textbf{Form} $\sqrt{a^2 - x^2}$: Let $x = a \sin \theta$, $dx = a \cos \theta \, d\theta$ \\
    (Relies on $1 - \sin^2 \theta = \cos^2 \theta$)
    \item \textbf{Form} $\sqrt{a^2 + x^2}$: Let $x = a \tan \theta$, $dx = a \sec^2 \theta \, d\theta$ \\
    (Relies on $1 + \tan^2 \theta = \sec^2 \theta$)
    \item \textbf{Form} $\sqrt{x^2 - a^2}$: Let $x = a \sec \theta$, $dx = a \sec \theta \tan \theta \, d\theta$ \\
    (Relies on $\sec^2 \theta - 1 = \tan^2 \theta$)
\end{itemize}
\end{multicols*}
\end{document}